\documentclass[11pt]{article}
\usepackage[T1]{fontenc}
\usepackage{lmodern}
\usepackage{amsmath,amssymb,amsthm}
\usepackage[margin=1.02in]{geometry}
\usepackage[hidelinks]{hyperref}

% Optional author details are separate from the project contact below.
\newcommand*{\AuthorName}{}
\newcommand*{\AuthorEmail}{}
\newcommand{\LiTwo}{\operatorname{Li}_2}
\newcommand{\pos}[1]{(#1)_{+}}
\newtheorem{theorem}{Theorem}
\newtheorem{proposition}[theorem]{Proposition}
\newtheorem{corollary}[theorem]{Corollary}

\title{The signed contribution of two- and three-prime products\\
to a truncated von Mangoldt convolution}
\author{\small \href{mailto:admin@cosmosgeek.org}{admin@cosmosgeek.org}\qquad \href{mailto:b321014@iiit-bh.ac.in}{b321014@iiit-bh.ac.in}}
\date{}

\begin{document}
\maketitle
\begin{abstract}
Consider the part of the divisor sum
\(\Lambda(n)=\sum_{d\mid n}\mu(d)\log(n/d)\) with \(d>X^t\), on
\(X<n\le 2X\). For \(0<t\le 1/3\) it is negative on every squarefree
product of two primes. On a product of three distinct primes \(r<s<q\)
it is positive exactly when \(s\le X^t<rs\), in which case its value
is \(\log q\). We evaluate the mass of this positive contribution,
also with a lower cutoff \(r>X^\delta\). Without the lower cutoff the
mass is \((\pi^2/12+o(1))X\), independently of \(t\). Even after
discarding all the negative three-prime contributions, the two-prime
terms outweigh the positive three-prime terms by an absolute multiple
of \(X\) for every fixed \(0\le\delta\le t\). The estimates concern
averages over the integers and over fixed reduced residue classes.
\end{abstract}

\section{The divisor weight and its signs}

Let \(\mu\) denote the M\"obius function, \(\Lambda\) the von Mangoldt
function and \(P^-(n)\) the least prime factor of \(n>1\). Fix \(X>1\)
and \(D=X^t\), where \(0<t\le 1/3\). For \(X<n\le 2X\) put
\begin{equation}\label{eq:weight}
 F_D(n)=\sum_{\substack{d\mid n\\ d>D}}\mu(d)\log(n/d)
       =\sum_{\substack{k\mid n\\1<k<n/D}}\mu(n/k)\log k.
\end{equation}
The omitted \(k=1\) term has weight zero. All logarithms are natural.
On composite squarefree \(n\) the identity
\(\sum_{d\mid n}\mu(d)\log(n/d)=\Lambda(n)=0\) gives
\begin{equation}\label{eq:complement}
 F_D(n)=-\sum_{\substack{d\mid n\\d\le D}}\mu(d)\log(n/d).
\end{equation}
For comparison, the usual truncated divisor sum in
Goldston--Y{\i}ld{\i}r{\i}m \cite{GY} uses \(\log(D/d)\).
On the right of \eqref{eq:complement}, replacing \(\log(n/d)\) by
\(\log(D/d)\) removes the term
\(-\log(n/D)\sum_{d\mid n,\,d\le D}\mu(d)\).

\begin{proposition}\label{prop:signs}
For primes \(r<s\) with \(X<rs\le 2X\),
\[
 F_D(rs)=
 \begin{cases}
 -\log(rs),&D<r,\\
 -\log r,&r\le D.
 \end{cases}
\]
For primes \(r<s<q\) with \(X<rsq\le 2X\),
\[
 F_D(rsq)=
 \begin{cases}
 -\log(rsq),&D<r,\\
 -\log r,&r\le D<s,\\
 0,&s\le D,\ rs\le D,\\
 \log q,&s\le D<rs.
 \end{cases}
\]
\end{proposition}
\begin{proof}
Since \(D^3\le X<n\), the largest prime exceeds \(D\) in each case.
For \(n=rs\), the possible divisors in \eqref{eq:complement} are
\(1\) and \(r\). For \(n=rsq\), they are \(1,r,s,rs\), subject to
the cutoff. If both \(r,s\) are included, the first three terms
give \(-\log(rsq)+\log(sq)+\log(rq)=\log q\). Including \(rs\)
subtracts \(\log q\). In the other cases the term \(d=1\) alone gives
\(-\log n\), and including only \(r\) in addition gives \(-\log r\).
\end{proof}

\section{The positive three-prime contribution}

Fix \(0\le\delta\le t\). Write \(\mathcal S_j(X,\delta)\) for the
squarefree integers in \((X,2X]\) having exactly \(j\) prime factors
and satisfying \(P^-(n)>X^\delta\). Thus \(\delta=0\) imposes no
additional restriction. Write \(v_+=\max(v,0)\).

\begin{theorem}\label{thm:positive}
As \(X\to\infty\) with \(t,\delta\) fixed,
\[
 \sum_{n\in\mathcal S_3(X,\delta)}\pos{F_D(n)}
   =\bigl(I(\delta/t)+o(1)\bigr)X .
\]
With \(\LiTwo(z)=-\int_0^z\log(1-y)\,dy/y\), the constant is
\begin{equation}\label{eq:constant}
 I(a)=
 \begin{cases}
 \pi^2/12, & a=0,\\[2pt]
 \dfrac12\left(\LiTwo(1-a)-\LiTwo(a)
 +\log(1/a)\log(1/(1-a))\right), &0<a<1/2,\\[3pt]
 \dfrac12\log^2(1/a),&1/2\le a\le 1 .
 \end{cases}
\end{equation}
\end{theorem}
\begin{proof}
By Proposition~\ref{prop:signs}, fix primes \(X^\delta<r<s\le D\)
with \(rs>D\). Then the possible primes \(q\) form the interval
\[
                  X/(rs)<q\le 2X/(rs).
\]
Every such \(q\) exceeds \(X^{1-2t}\) and hence \(D\). With
\(\vartheta(y)=\sum_{q\le y}\log q\), the prime number theorem gives
\(\vartheta(2y)-\vartheta(y)=(1+o(1))y\), uniformly for
\(y\ge X^{1-2t}\). It follows that the sum under consideration equals
\begin{equation}\label{eq:harm}
  (1+o(1))X
  \sum_{\substack{X^\delta<r<s\le D\\rs>D}}\frac{1}{rs}.
\end{equation}
Here and below variables denoting prime factors in harmonic sums run
over primes.

For fixed \(0<A<B\), partial summation of the prime number theorem
gives \(\sum_{X^A<r\le X^B}1/r\to\log(B/A)\). Thus the measure
putting mass \(1/r\) at \(u=\log r/\log X\) converges on compact
subintervals of \((0,\infty)\) to \(du/u\). For \(\delta>0\) this
applies to the product measure in \((u,v)\). The lines \(u=v\) and
\(u+v=t\), and the endpoints of the box, have measure zero in the
limit. Symmetry therefore gives the limit of the sum in \eqref{eq:harm}
as
\begin{equation}\label{eq:double}
 \frac12\int_{\delta}^{t}\int_{\max(\delta,t-u)}^{t}
                                     \frac{dv\,du}{uv}.
\end{equation}
The factor \(1/2\) counts each unordered pair of distinct small primes
once.

The limit at \(\delta=0\) requires control near \(u=0\).
For \(0<\epsilon<t/2\) and \(r\le X^\epsilon\),
the interval \(D/r<s\le D\) is covered by \(O(\log(2r))\) dyadic
intervals with endpoints \(\gg X^{t/2}\). The prime counting upper
bound \(\pi(y)\ll y/\log y\) on each interval gives
\[
 \sum_{r\le X^\epsilon}\frac1r
   \sum_{D/r<s\le D}\frac1s
 \ll_t \frac1{\log X}\sum_{r\le X^\epsilon}\frac{\log(2r)}r
 \ll_t \epsilon+\frac1{\log X}.
\]
The last bound follows by partial summation. We may first fix
\(\epsilon\), take \(X\to\infty\), and then let \(\epsilon\downarrow0\).

Scale \(u,v\) in \eqref{eq:double} by \(t\). For \(0<a<1/2\),
twice the integral is
\[
 \int_a^{1-a}\frac{-\log(1-y)}{y}\,dy
  +\log(1/a)\int_{1-a}^1\frac{dy}{y}.
\]
This gives the middle line of \eqref{eq:constant}. At \(a=0\)
it gives \(\LiTwo(1)=\sum_{j\ge1}j^{-2}=\pi^2/6\) before division
by two. At \(a\ge1/2\), the entire open square \((a,1)^2\)
satisfies \(u/t+v/t>1\), which gives the last line.
\end{proof}

\section{The signed mean}

\begin{theorem}\label{thm:signed}
For \(0<t\le 1/3\) and \(0\le\delta\le t\), put
\begin{align}
 B&=\log((1-\delta)/t),\nonumber\\
 J_0&=\int_t^{1/3}\int_u^{(1-u)/2}
                    \frac{dv\,du}{uv(1-u-v)},\label{eq:Js}\\
 J_1&=\int_\delta^t\int_t^{(1-u)/2}
                    \frac{dv\,du}{v(1-u-v)}.\nonumber
\end{align}
For fixed \(t,\delta\), as \(X\to\infty\),
\begin{align*}
 X^{-1}\!\sum_{n\in\mathcal S_2(X,\delta)}F_D(n)&=-B+o(1),\\
 X^{-1}\!\sum_{n\in\mathcal S_3(X,\delta)}F_D(n)
         &=I(\delta/t)-J_0-J_1+o(1).
\end{align*}
Even if all negative three-prime terms are discarded, one has
\[
 \lim_{X\to\infty}\frac1X\left(
   \sum_{n\in\mathcal S_2(X,\delta)}F_D(n)+
   \sum_{n\in\mathcal S_3(X,\delta)}\pos{F_D(n)}\right)
 \le -\left(\log(5/2)-\frac{\pi^2}{12}\right)<0.
\]
\end{theorem}
\begin{proof}
Write \(L=\log X\). If \(X^\delta<r\le D\), the contribution of \(rs\)
to the negative mass is \(\log r\) for each prime
\(X/r<s\le 2X/r\). If \(D<r\le\sqrt X\), the negative mass is
\(\log(rs)=L+O(1)\) for each such prime. The range
\(\sqrt X<r\le\sqrt{2X}\) contributes \(O(X/L)\), using
\(\pi(2X/r)\ll X/(rL)\) and
\(\sum_{\sqrt X<r\le\sqrt{2X}}1/r\ll 1/L\).
The prime number theorem gives
\(\pi(2y)-\pi(y)=(1+o(1))y/\log y\) uniformly at \(y=X/r\) in
the first two ranges. For the first range we use the measure with
mass \(\log r/(rL)\) at \(u=\log r/L\); it converges to \(du\) on
\([0,t]\) by partial summation, also at the endpoint \(0\).
For the second range we use the prime harmonic measure \(du/u\)
from the proof of Theorem~\ref{thm:positive}. The two limiting negative
masses divided by \(X\) consequently add to
\[
  \int_\delta^t\frac{du}{1-u}+
  \int_t^{1/2}\frac{du}{u(1-u)}=\log((1-\delta)/t)=B .
\]

For three primes \(r<s<q\), the negative mass in
Proposition~\ref{prop:signs} is \(\log(rsq)\) when \(r>D\), and
\(\log r\) when \(X^\delta<r\le D<s\). Set \(u=\log r/L\),
\(v=\log s/L\). If \(rs^2\le X\), the order \(q>s\) holds for
every \(X/(rs)<q\le 2X/(rs)\). If \(rs^2>2X\), there is no
contribution. In the intermediate strip
\(1<u+2v\le 1+\log(2)/L\), the number of choices is at most
\(\pi(2X/(rs))\ll X/(rsL)\), since \(s>X^t\) and
\(X/(rs)\ge s/2\).
After division by \(X\) the strip has vanishing mass: the relevant
measures on the box have limits \(du\,dv/(uv)\) when \(u\ge t\)
and \(du\,dv/v\) when \(0\le u\le t\). Each assigns measure zero
to \(u+2v=1\). This also justifies neglecting the ordering boundary
\(u=v\).

On \(u+2v\le1\) we have \(\log(X/(rs))=(1-u-v)L\ge tL\).
Using the prime number theorem for \(q\) now gives the integral
\(J_0\) for the first negative case, since \(\log(rsq)=L+O(1)\).
In the second case \(\log r=uL\) cancels the \(u\) in the harmonic
density, giving \(J_1\). Theorem~\ref{thm:positive} supplies the
positive term.

Finally, if \(\delta\le t/2\), then \(B\ge\log(5/2)\) and
\(I(\delta/t)\le\pi^2/12\). If \(t/2\le\delta\le t\), then
\(B\ge\log 2\) and \(I(\delta/t)\le(\log 2)^2/2\). The latter
gap exceeds \(\log(5/2)-\pi^2/12\). The little-oh errors above
are for fixed parameters.
\end{proof}

\section{Fixed congruence restrictions and shift two}

\begin{corollary}\label{cor:local}
Restrict any of the sums in Theorems~\ref{thm:positive}
and~\ref{thm:signed} to \(n\equiv b\pmod m\), where \(m\) is fixed
and \((b,m)=1\). Each leading constant is divided by \(\varphi(m)\).
In particular, requiring \((n+2,m)=1\) multiplies each leading
constant by
\[
       \prod_{\substack{\ell\mid m\\\ell>2\ {\rm prime}}}
                     \frac{\ell-2}{\ell-1}.
\]
\end{corollary}
\begin{proof}
For \(\delta>0\), all prime factors of \(n\) exceed the prime divisors
of \(m\) once \(X\) is large. Once the smaller factors are fixed, the
largest prime must belong to one reduced class modulo \(m\).
The prime number theorem for a fixed arithmetic progression replaces
each use of the prime number theorem in its interval by a factor
\(1/\varphi(m)\); see \cite{Tao}. For \(\delta=0\) remove \(r\le
X^\epsilon\) first. The harmonic tail estimate in
Theorem~\ref{thm:positive} and the integrals for \(B,J_1\) in
Theorem~\ref{thm:signed} bound its absolute weight by
\(O_t(\epsilon X)+o(X)\). Thus nonunit residue classes contribute
\(o(X)\). Among the reduced residue classes, the additional
condition excludes exactly one unit class modulo each odd
\(\ell\mid m\) and excludes none modulo \(2\). The Chinese
remainder theorem gives the product.
\end{proof}

For example, at \(t=0.24995\) and \(\delta=1/20\),
\[
  I(\delta/t)=0.6114186198\ldots,\qquad
  B=1.3352010867\ldots .
\]
Thus the signed sum with negative three-prime terms omitted is
\((-0.7237824669\ldots+o(1))X\). Including them gives
\((-1.0754928449\ldots+o(1))X\) by \eqref{eq:Js}.
Requiring \(3\nmid n+2\) divides either main term by two.
A condition that \(n+2\) itself be
prime requires estimates uniform in growing arithmetic restrictions
or a direct estimate for the two integers together. A fixed-modulus
mean leaves that correlation undetermined.

\begin{thebibliography}{9}
\bibitem{GY} D.~A. Goldston and C.~Y. Y{\i}ld{\i}r{\i}m,
  \emph{Higher correlations of divisor sums related to primes I:
  triple correlations}, Integers \textbf{3} (2003), A05.
  \href{https://arxiv.org/abs/math/0111212}{arXiv:math/0111212}.
\bibitem{Tao} T.~Tao, \emph{The prime number theorem in arithmetic
  progressions, and dueling conspiracies}, 24 September 2009.
  \href{https://terrytao.wordpress.com/2009/09/24/the-prime-number-theorem-in-arithmetic-progressions-and-dueling-conspiracies/}{Author's notes}.
\end{thebibliography}
\end{document}
